\section{Introduction}
\label{sec:intro}

Garment manipulation is a contact-rich, bimanual task that remains difficult
to automate directly; learning-based approaches (e.g.\ vision-language-action
policies such as $\pi_{0.5}$) depend on high-quality demonstration data,
which in turn depends on a teleoperation interface that is precise,
low-latency, and comfortable enough for long collection sessions. This paper
documents the teleoperation stack of the \texttt{so101\_garment} rig: two
SO-101 arms (the LeRobot/SO-ARM100 lineage~\citep{soarm100}) mounted on a 3D-printed
platform, teleoperated with Meta Quest hand controllers, with a MuJoCo
digital twin used both for operator rehearsal and for quantitative
benchmarking of design alternatives.

Three properties of the SO-101 shape every design decision. First, each arm
has only five degrees of freedom (shoulder pan, shoulder lift, elbow flex,
wrist flex, wrist roll): there is no wrist-yaw joint, so full 6D pose
tracking is structurally over-constrained~--- a naive 6D tracker trades
position error against an orientation error it can never eliminate. Second,
the arm is small (\(\approx\)\SI{0.41}{m} maximal reach from the shoulder-lift
pivot) while a human operator's natural gestures easily exceed that reach, so
out-of-envelope hand targets are the norm, not the exception. Third, the
STS3215 bus servos are position-controlled with modest stiffness, so command
smoothness directly affects hardware wear and tracking fidelity.

This document makes five contributions:
\begin{enumerate}
  \item a description of the production retargeting pipeline: grip-clutched
        absolute-position mapping with One-Euro filtering and the
        \emph{armplane} orientation construction, which builds only
        orientations the 5-DoF arm can actually realize
        (Section~\ref{sec:method});
  \item an analytic \emph{workspace envelope} for the SO-101 with four
        selectable out-of-envelope target policies~--- warn, project,
        freeze, and slowdown~--- applied uniformly to every teleoperation
        input source (Section~\ref{sec:envelope});
  \item the integration of a \emph{Telegrip-style split IK}
        \citep{telegrip} as a pluggable method: position-only damped least
        squares on the three proximal joints with a direct analytic wrist
        mapping, targeting wrist agility (Section~\ref{sec:telegrip});
  \item a simulation benchmark comparing six IK methods on tracking,
        wrist-agility, and out-of-envelope suites, with metrics that
        include orientation error and roll phase lag
        (Section~\ref{sec:experiments});
  \item a within-subjects \emph{user study protocol} (\(n\approx 5\))
        pairing the objective benchmark with subjective evaluation~---
        NASA-TLX workload, system usability, custom teleoperation-feel
        scales, and semi-structured interviews~--- on a bimanual
        coordination task (Section~\ref{sec:exp-userstudy}).
\end{enumerate}

The paper is a \emph{living document}: it is the canonical reference for the
teleoperation side of the repository and is updated in the same branch as
any change to the methods, mappings, envelope, or benchmark results.
